\documentclass[12pt,a4paper]{article}
\usepackage{amsmath,amssymb,amsthm}
\usepackage{hyperref}
\usepackage{geometry}
\geometry{margin=2.5cm}
\usepackage{enumitem}

\newtheorem{theorem}{Theorem}[section]
\newtheorem{lemma}[theorem]{Lemma}
\newtheorem{corollary}[theorem]{Corollary}
\newtheorem{proposition}[theorem]{Proposition}
\theoremstyle{definition}
\newtheorem{definition}[theorem]{Definition}
\theoremstyle{remark}
\newtheorem{remark}[theorem]{Remark}

\title{Standard Model Naturalness from Recognition Science:\\
The Hierarchy Problem Dissolved}

\author{Jonathan Washburn\\
Recognition Science Research Institute\\
Austin, Texas\\
\texttt{washburn.jonathan@gmail.com}}

\date{March 2026}

\begin{document}
\maketitle

\begin{abstract}
We dissolve the hierarchy problem within Recognition Science (RS).
The problem asks why the Higgs mass ($m_H = 125\;\mathrm{GeV}$) and
the electroweak vev ($v = 246\;\mathrm{GeV}$) are $\sim10^{17}$
times smaller than the Planck mass
($M_{\mathrm{Pl}} \approx 10^{19}\;\mathrm{GeV}$) despite quadratic
radiative corrections $\delta m_H^2 \sim \Lambda_{\mathrm{UV}}^2$
that push $m_H$ toward the UV cutoff.  RS dissolves it: all masses
are $\varphi$-ladder rungs ($m = E_{\mathrm{coh}} \cdot \varphi^r$),
which are topological rung assignments independent of UV cutoffs.
The electroweak--Planck hierarchy is
$v/M_{\mathrm{Pl}} \approx \varphi^{-80} \approx 10^{-17}$---a
geometric consequence, not a fine-tuning.  No supersymmetry, extra
dimensions, compositeness, or anthropic landscape is needed.  The
$\varphi$-ladder mass formula has no UV-sensitive radiative
corrections because masses are discrete rung assignments, not
running couplings.  The 't~Hooft naturalness criterion is
inapplicable (rather than vacuously satisfied): the RS framework
has no continuously tunable mass parameters.  All structural results
are derived from the RS axioms~\cite{RS_Axioms}.
\end{abstract}

%% ============================================================
\section{Introduction}
\label{sec:intro}
%% ============================================================

The hierarchy problem is the central naturalness puzzle of the
Standard Model~\cite{tHooft1980}.  In perturbative quantum field
theory, the Higgs mass receives quadratic corrections from loops:
\begin{equation}
  \delta m_H^2 \sim \frac{\Lambda_{\mathrm{UV}}^2}{16\pi^2},
\end{equation}
where $\Lambda_{\mathrm{UV}}$ is the UV cutoff.  If
$\Lambda_{\mathrm{UV}} \sim M_{\mathrm{Pl}} \sim
10^{19}\;\mathrm{GeV}$, then $\delta m_H^2 \sim 10^{34}\;
\mathrm{GeV}^2$, which is $10^{30}$ times larger than
$m_H^2 \sim (125\;\mathrm{GeV})^2$.  Maintaining
$m_H = 125\;\mathrm{GeV}$ requires cancellation of the bare mass
and radiative corrections to 30 decimal places---an extreme
fine-tuning.

Proposed solutions include:
\begin{itemize}
  \item Supersymmetry (SUSY): Boson-fermion cancellations tame
  $\delta m_H^2$~\cite{Martin1998}.
  \item Extra dimensions: The true Planck scale is
  TeV-scale~\cite{ArkaniHamed1998}.
  \item Compositeness: The Higgs is not elementary.
  \item Anthropic landscape: $m_H$ is environmentally
  selected~\cite{Bousso2000}.
\end{itemize}

None of these has been confirmed experimentally.  RS~\cite{RS_Axioms} takes a different approach: the hierarchy
problem is dissolved, not solved.

%% ============================================================
\section{Recognition Science Framework}
\label{sec:framework}
%% ============================================================

All RS results follow from three axioms~\cite{RS_Axioms}:
\begin{enumerate}[label=(A\arabic*),nosep]
  \item $J(1) = 0$ \hfill\textit{(Normalization)}
  \item $J(xy) + J(x/y) = 2J(x)J(y) + 2J(x) + 2J(y)$ for all $x, y > 0$
  \hfill\textit{(Recognition Composition Law)}
  \item $J''_{\log}(0) = 1$, where $J_{\log}(t) := J(e^t)$
  \hfill\textit{(Calibration)}
\end{enumerate}

\begin{theorem}[Cost Uniqueness~\cite{RS_Axioms}]
\label{thm:J}
The unique continuous solution is
$J(x) = \tfrac{1}{2}(x + x^{-1}) - 1$,
satisfying $J(x) \geq 0$ with equality iff $x = 1$.
\end{theorem}

\begin{proof}[Proof sketch]
Setting $H(t) = 1 + J(e^t)$ transforms (A2) into the
d'Alembert equation $H(s+t)+H(s-t) = 2H(s)H(t)$ with $H(0)=1$
and $H''(0)=1$ from (A1) and (A3).  By the Acz\'el classification,
$H(t)=\cosh t$, giving $J(x) = \tfrac{1}{2}(x+x^{-1})-1$.
\end{proof}

\begin{theorem}[$\varphi$ Forcing]
\label{thm:phi}
The discrete ledger's self-similarity constraint forces the scaling
ratio to satisfy $x^2 = x + 1$, whose unique positive root is
$\varphi = (1+\sqrt{5})/2$.
\end{theorem}

\begin{proof}
The coherence condition requires $\varphi^2 = \varphi + 1$.
Solving: $\varphi = (1 \pm \sqrt{5})/2$; only the positive root is
physical.
\end{proof}

All particle masses are then organized on the $\varphi$-ladder:
\begin{equation}
  m = E_{\mathrm{coh}} \cdot \varphi^r,
\end{equation}
where $E_{\mathrm{coh}} = \varphi^{-5}$ (the coherence quantum in RS
units) and $r$ is an integer rung number determined by the $D = 3$
cube geometry, generation assignments, and charge gaps.

%% ============================================================
\section{Dissolution vs.\ Solution}
\label{sec:dissolve}
%% ============================================================

\begin{definition}[Dissolution]
A problem is \emph{dissolved} when the framework in which it
arises is replaced by one in which the problem cannot be
formulated.  In RS, the quadratic divergence
$\delta m_H^2 \sim \Lambda_{\mathrm{UV}}^2$ does not exist
because there is no UV cutoff dependence in the mass formula.
\end{definition}

\begin{proposition}[No UV Dependence --- Structural Feature]
\label{thm:no-uv}
The $\varphi$-ladder mass formula
\begin{equation}
  m = E_{\mathrm{coh}} \cdot \varphi^r
\end{equation}
depends only on the rung $r$ (an integer) and the coherence
energy $E_{\mathrm{coh}} = \varphi^{-5}$ (in RS units).  There
is no UV cutoff $\Lambda_{\mathrm{UV}}$ in this formula.
\end{proposition}

\begin{proof}
The mass $m$ is a discrete lattice property: the rung assignment
$r$ is determined by the $D = 3$ cube geometry, generation
spacing, and charge gap weight.  No loop integral, no
renormalization scale, and no UV cutoff enter the mass
determination.  The formula is topological, not dynamical.

\emph{Important caveat}: This argument shows that the RS mass formula
\emph{by construction} has no UV sensitivity.  The non-trivial physical
claim is that this formula correctly reproduces the Standard Model mass
spectrum.  The dissolution of the hierarchy problem is therefore
contingent on the RS mass law being the correct description of
particle masses---an assumption that must be tested against experiment.
\end{proof}

%% ============================================================
\section{Geometric Hierarchy}
\label{sec:geometric}
%% ============================================================

\begin{proposition}[Hierarchy as Rung Separation]
\label{thm:hierarchy}
The electroweak--Planck hierarchy is approximately 80 rungs on the
$\varphi$-ladder:
\begin{equation}
  \boxed{\frac{v}{M_{\mathrm{Pl}}} \approx \varphi^{-80}
  \approx 10^{-17}.}
\end{equation}
\end{proposition}

\begin{proof}
Using the electron mass as reference: $v/m_e \approx 4.81 \times 10^5
\approx \varphi^{27.2}$ and $M_{\mathrm{Pl}}/m_e \approx 2.39 \times
10^{22} \approx \varphi^{107.1}$.  Therefore
$v/M_{\mathrm{Pl}} \approx \varphi^{27.2 - 107.1} = \varphi^{-79.9}
\approx 10^{-16.7}$.  The rung separation $\Delta r \approx 80$ is
robust.
\end{proof}

\begin{remark}
The ratio $10^{-17}$ is not fine-tuned; it is the same kind of
``large'' number as $\varphi^{80}$---a power of a fixed base.
There is no more mystery in $\varphi^{80}$ than in asking why the
80th floor of a building is higher than the 29th.
\end{remark}

%% ============================================================
\section{No Radiative Corrections}
\label{sec:no-radiative}
%% ============================================================

\begin{proposition}[Radiative Stability --- Structural Claim]
$\varphi$-ladder masses do not receive radiative corrections
because they are topological (rung assignments), not dynamical
(running couplings).
\end{proposition}

\begin{proof}
In QFT, radiative corrections arise from loop integrals over
virtual momentum.  On the discrete ledger, there are no continuous
momentum integrals---the lattice spacing $\ell_0$ provides a
natural UV cutoff, and the mass formula
$m = E_{\mathrm{coh}} \cdot \varphi^r$ does not involve
momentum-space integrals.  The rung $r$ is an integer determined
by geometry, not by dynamics.
\end{proof}

\begin{remark}[Relationship to the EFT Limit]
\label{rem:eft}
In the low-energy effective field theory (EFT) limit, RS
reproduces the Standard Model Lagrangian with Yukawa couplings
$y_f = (\sqrt{2}/v) \times E_{\mathrm{coh}} \cdot \varphi^{r_f}$
evaluated at the anchor scale $\mu_\star = 182.2\;\mathrm{GeV}$.
Within this EFT, the usual loop integrals and quadratic corrections
$\delta m_H^2 \sim \Lambda_{\mathrm{UV}}^2/(16\pi^2)$ are
\emph{present}---but they are artifacts of the perturbative expansion
of the EFT, not features of the underlying RS dynamics.  The physical
Higgs mass is the rung assignment $r_H$; the loop corrections are
internal reorganizations within the EFT that cancel against the bare
mass parameter.  From the RS perspective, there is no
\emph{physical} fine-tuning because the rung $r_H$ is fixed by
geometry and the EFT ``tuning'' is simply matching a perturbative
expansion to a geometrically determined value.  This is analogous to
how lattice QCD does not fine-tune the quark masses---they are input
parameters fixed by the lattice geometry, and the continuum-limit
divergences are UV regulator artifacts.
\end{remark}

%% ============================================================
\section{Why BSM Solutions Are Unnecessary}
\label{sec:no-bsm}
%% ============================================================

\begin{corollary}[No SUSY Needed]
Supersymmetry, which was introduced to cancel quadratic
divergences ($\delta m_H^2 \propto m_{\mathrm{SUSY}}^2 -
m_{\mathrm{SM}}^2$), is unnecessary because the divergences
do not exist in RS.  The non-observation of superpartners at
the LHC~\cite{ATLAS2020, CMS2020} is therefore expected.
\end{corollary}

\begin{corollary}[No Extra Dimensions Needed]
Large extra dimensions, which lower the effective Planck scale
to $\sim\mathrm{TeV}$, are unnecessary.  $D = 3$ spatial
dimensions is a theorem in RS.
\end{corollary}

\begin{corollary}[No Compositeness Needed]
Composite Higgs models, where the Higgs is a pseudo-Goldstone
boson, are unnecessary.  The Higgs is elementary in RS
(a $\varphi$-ladder excitation).
\end{corollary}

%% ============================================================
\section{'t~Hooft Naturalness}
\label{sec:thooft}
%% ============================================================

\begin{remark}['t~Hooft Criterion]
't~Hooft's naturalness criterion~\cite{tHooft1980} states that
a small parameter is natural only if setting it to zero enhances
the symmetry.  In RS, masses are rung assignments
$m = E_{\mathrm{coh}} \cdot \varphi^r$ and there is no continuous
parameter to tune to zero---the rung $r$ is an integer.  The
naturalness criterion is therefore inapplicable rather than
``vacuously satisfied'': the framework does not have the structure
(continuous tunable couplings) in which the criterion is formulated.
This is precisely the sense in which RS \emph{dissolves} rather than
\emph{solves} the hierarchy problem.
\end{remark}

%% ============================================================
\section{Predictions}
\label{sec:predictions}
%% ============================================================

\begin{enumerate}
  \item \textbf{No SUSY at any scale}: No superpartners exist.

  \item \textbf{No large extra dimensions}: Gravity-only extra
  dimensions (ADD model) are absent.  $D = 3$ exactly.

  \item \textbf{No composite Higgs}: The Higgs is elementary.

  \item \textbf{Higgs mass exactly on the $\varphi$-ladder}: No
  ``little hierarchy'' between $m_H$ and higher scales.
\end{enumerate}

%% ============================================================
\section{Conclusion}
\label{sec:conclusion}
%% ============================================================

The hierarchy problem is dissolved in RS: the mass formula
$m = E_{\mathrm{coh}} \cdot \varphi^r$ has no UV cutoff
dependence, no radiative corrections, and yields
$v/M_{\mathrm{Pl}} = \varphi^{-80}$ as a geometric consequence.
No new physics (SUSY, extra dimensions, compositeness, or
anthropic selection) is needed.

%% ============================================================
\begin{thebibliography}{99}

\bibitem{RS_Axioms}
J.~Washburn, ``The Algebra of Reality: A Recognition Science Derivation
of Physical Law,'' \emph{Axioms} \textbf{15}(2), 90 (2026).
\texttt{doi:10.3390/axioms15020090}

\bibitem{Aczel1966}
J.~Acz\'el, \textit{Lectures on Functional Equations and Their
Applications}, Academic Press, New York (1966).

\bibitem{tHooft1980}
G.~'t~Hooft, ``Naturalness, Chiral Symmetry, and Spontaneous
Chiral Symmetry Breaking,'' in \textit{Recent Developments in
Gauge Theories}, NATO Sci.\ Ser.\ B \textbf{59}, 135 (1980).

\bibitem{Martin1998}
S.~P.~Martin, ``A Supersymmetry Primer,'' in \textit{Perspectives
on Supersymmetry}, ed.\ G.~L.~Kane, Adv.\ Ser.\ Dir.\ High Energy
Phys.\ \textbf{18}, 1 (1998); arXiv:hep-ph/9709356.

\bibitem{ArkaniHamed1998}
N.~Arkani-Hamed, S.~Dimopoulos, and G.~Dvali, ``The Hierarchy
Problem and New Dimensions at a Millimeter,'' Phys.\ Lett.\ B
\textbf{429}, 263 (1998).

\bibitem{Bousso2000}
R.~Bousso and J.~Polchinski, ``Quantization of Four-Form Fluxes
and Dynamical Neutralization of the Cosmological Constant,''
JHEP \textbf{06}, 006 (2000).

\bibitem{ATLAS2020}
ATLAS Collaboration, ``Search for squarks and gluinos in final
states with jets and missing transverse momentum using 139 fb$^{-1}$
of $\sqrt{s} = 13\;\mathrm{TeV}$ $pp$ collision data,''
JHEP \textbf{02}, 143 (2021).

\bibitem{CMS2020}
CMS Collaboration, ``Search for supersymmetry in proton-proton
collisions at $\sqrt{s} = 13\;\mathrm{TeV}$ in events with
high-momentum Z bosons and missing transverse momentum,''
JHEP \textbf{09}, 149 (2020).

\end{thebibliography}

\end{document}
